% ── Interlude 01: The Surface of Inscription ─────────────────
% Object class: manuscript, encoded surface
% Appears after Part I (Guyon) before Part II (Mackandal)
% No named characters. No date. No location.

\begin{interlude}{The Surface of Inscription}

A surface receives marks.

This is its primary function, though it was not made for this. It
was made from something else---skin, pulp, reed, clay---and the
making involved pressure, heat, the removal of what was not needed.
What remains is a capacity: to hold an impression without
immediately losing it.

The marks vary. Some are deep, cut into the material by force.
Others sit on the surface, almost weightless, ready to smear at
the first dampness. Between these extremes lies most of what
has ever been written: marks that seem permanent until examined
closely, at which point they reveal themselves as merely stable,
held in place by conditions that could, in principle, change.

A surface with marks on it is not yet a text.

Something else is required---a reader, a code, a set of
transformations that convert the pattern of marks into something
that can move from one mind to another without losing its
structure. Without this, the marks remain what they are: local
disturbances in a material, persisting because nothing has yet
disturbed them further.

The same surface, read by different systems, yields different
texts. This is not a failure of the surface. It is a property of
encoding: meaning is not stored in the marks themselves, but in
the relation between the marks and the process that interprets
them. Change the process, and the text changes, even if the
surface does not.

\medskip

Some surfaces are written on in layers.

The earlier marks are not fully erased before the later ones are
added. They persist beneath, faintly, recoverable under certain
light. The text that appears on the surface is not the only text
the surface contains.

This is sometimes a problem and sometimes a resource.

When it is a problem, it is because the earlier marks interfere
with the reading of the later ones, introducing ambiguity where
none was intended. When it is a resource, it is because the
earlier marks carry something the later ones cannot: a record of
what the surface held before it held this, evidence of a prior
configuration that the current text has not entirely displaced.

A surface that has been written on many times remembers, imperfectly,
everything it has held.

\medskip

There are surfaces that were never meant to be read.

They were written on in order to fix something in place---to
prevent it from shifting, to hold it against the pressure of what
surrounds it. The marks are not addressed to a reader. They are
addressed to the material itself, or to whatever it is that makes
materials hold their shape.

Such a surface, found later, presents a difficulty.

It can be read, in the sense that the marks can be interpreted
according to whatever system the finder brings to it. But this
reading may have nothing to do with why the marks were made. The
surface is legible, but its legibility is a secondary property,
not its function.

Most surfaces are like this, if examined closely enough.

They were made for reasons that reading cannot fully recover. The
marks mean something, but not only what they say.

\medskip

The surface persists.

The marks on it persist, or most of them. The conditions that
produced them have long since changed---the hand that made them,
the light in the room, the question that required an answer---but
the marks themselves remain, holding their positions relative to
one another, maintaining the pattern that was fixed into the
material at the moment of inscription.

Whether this constitutes memory is a question the surface cannot
answer.

It holds what it was given.

It offers this, without interpretation, to whatever comes next.

\end{interlude}
